% Cote 119: offprints, reports, courses, CV and letters (typed and printed part): archive cote 119, site page 207 (PDF page 208).
% typescript. Draft compared: none.
% AI-transcribed by Claude Opus 5.5 (claude-opus-5-5, adaptive thinking, effort high), Message Batches API, batch r119_full (msgbatch_01DwZNm5Y7YsKBWyAN4sFs8w), 2026-11-13; not yet checked by eye.
\scanpage{119/207}
\mspage{- 12 -}

I. Manin [*] et M. Demazure [*]), mais dès qu'on veut aller plus loin et for-%
muler des propriétés fondamentales ``naturelles'', on bute sur des conjectures
actuellement indémontrables, comme celles de Weil ou de Tate, et d'autres ana-%
logues que la notion même de motif suggère irrésistiblement. Ces propriétés ont
fait l'objet de nombreuses conversations privées et de plusieurs exposés
publics, mais n'ont jamais fait l'objet d'une publication, puisqu'il n'est pas
d'usage en mathématique (contrairement à la physique) de publier un rêve, si
cohérent soit-il, et de suivre jusqu'au bout où ses divers éléments nous peuvent
entrainer. Il est évident pourtant, pour quiconque se plonge suffisamment dans
la cohomologie des variétés algébriques, `` qu'il y a quelque chose'' - que
``les motifs existent''. Il y a quelques années encore, j'ai joué avec l'idée
d'écrire contrairement à l'usage, un livre entièrement conjectural sur les
motifs - une sorte de science-fiction mathématique. J'en ai été empêché par
des tâches plus urgentes que des tâches de mathématicien, et je doute fort
actuellement qu'un tel livre soit jamais écrit, ni qu'on arrive jamais (même
conjecturalement) à se faire une idée d'ensemble à la fois précise et suffisam-%
ment vaste sur le formalisme des motifs. Avant qu'on n'y parvienne, il sera
sans doute devenu évident pour tous, sous la poussée des événements, la science
spéculative et parcellarisée ne faisant plus vivre son homme, qu'il est des
tâches plus urgentes que de mettre sur pied même la plus belle théorie du
monde, conjecturale ou non.

\mspage{\hand{207}}
